\answerpaper{数学 B 类}

\answerproblem{一、}{15}
三角形三条中线的交点称为三角形的重心．在四面体 $ABCD$ 中，记 $\overrightarrow{e_1}:=\overrightarrow{AB}$，$\overrightarrow{e_2}:=\overrightarrow{AC}$，$\overrightarrow{e_3}:=\overrightarrow{AD}$，并设 $O_1,O_2$ 和 $O_3$ 分别为 $\triangle BCD$，$\triangle ACD$ 和 $\triangle ABD$ 的重心．
\begin{enumerate}[label=(\arabic*)]
  \item 在坐标系 $\{A,\overrightarrow{e_1},\overrightarrow{e_2},\overrightarrow{e_3}\}$ 下，求 $O_1$ 点的坐标 $(x,y,z)$，其中 $\overrightarrow{AO_1}=x\overrightarrow{e_1}+y\overrightarrow{e_2}+z\overrightarrow{e_3}$；
  \item 证明三直线 $AO_1$，$BO_2$ 及 $CO_3$ 相交于一点 $P$；
  \item 在坐标系 $\{A,\overrightarrow{e_1},\overrightarrow{e_2},\overrightarrow{e_3}\}$ 下，求 (2) 中交点 $P$ 的坐标．
\end{enumerate}
\begin{solution}
\textbf{(1)} 由于 $O_1$ 是 $\triangle BCD$ 的重心，直线 $BO_1$ 与对边 $CD$ 交于点 $E$，则 $\dfrac{BO_1}{O_1E}=2$．从而
\[
  \overrightarrow{AO_1}=\frac{1}{3}\overrightarrow{AB}+\frac{2}{3}\overrightarrow{AE}
  =\frac{1}{3}\overrightarrow{AB}+\frac{1}{3}\left(\overrightarrow{AC}+\overrightarrow{AD}\right).
\]
从而在坐标系 $\{A,\overrightarrow{e_1},\overrightarrow{e_2},\overrightarrow{e_3}\}$ 下，$O_1$ 点的坐标为 $\left(\frac{1}{3},\frac{1}{3},\frac{1}{3}\right)$．

\textbf{(2)} 类似可得 $O_2,O_3$ 的坐标分别为 $\left(0,\frac{1}{3},\frac{1}{3}\right)$、$\left(\frac{1}{3},0,\frac{1}{3}\right)$．

设 $AO_1$ 与 $BO_2$ 的交点为 $P(x,y,z)$．令 $\overrightarrow{AP}=\lambda\overrightarrow{AO_1}$，$\overrightarrow{BP}=\mu\overrightarrow{BO_2}$．由 $\overrightarrow{AP}+\overrightarrow{PB}+\overrightarrow{BA}=\vec{0}$ 得：
\[
  \lambda\overrightarrow{AO_1}-\mu\overrightarrow{BO_2}-\overrightarrow{AB}=\vec{0}.
\]
即
\[
  \lambda\left(\frac{1}{3}\vec{e}_1+\frac{1}{3}\vec{e}_2+\frac{1}{3}\vec{e}_3\right)
  -\mu\left(\frac{1}{3}\vec{e}_2+\frac{1}{3}\vec{e}_3-\vec{e}_1\right)-\vec{e}_1=\vec{0}.
\]
解得 $\lambda=\mu=\dfrac{3}{4}$．故
\[
  \overrightarrow{AP}=\lambda\overrightarrow{AO_1}
  =\frac{3}{4}\left(\frac{1}{3}\vec{e}_1+\frac{1}{3}\vec{e}_2+\frac{1}{3}\vec{e}_3\right)
  =\frac{1}{4}(\vec{e}_1+\vec{e}_2+\vec{e}_3).
\]
即 $P$ 的坐标为 $\left(\frac{1}{4},\frac{1}{4},\frac{1}{4}\right)$．

此时
\[
  \overrightarrow{CP}=\overrightarrow{AP}-\overrightarrow{AC}
  =\frac{1}{4}(\vec{e}_1+\vec{e}_2+\vec{e}_3)-\vec{e}_2
  =\frac{1}{4}\vec{e}_1-\frac{3}{4}\vec{e}_2+\frac{1}{4}\vec{e}_3.
\]
又 $\overrightarrow{CP}=3\cdot\overrightarrow{PO_3}$，
\[
  \overrightarrow{PO_3}=\overrightarrow{AO_3}-\overrightarrow{AP}
  =\frac{1}{3}\vec{e}_1+\frac{1}{3}\vec{e}_3-\frac{1}{4}(\vec{e}_1+\vec{e}_2+\vec{e}_3)
  =\frac{1}{12}\vec{e}_1-\frac{1}{4}\vec{e}_2+\frac{1}{12}\vec{e}_3.
\]
故 $C,P,O_3$ 三点共线．从而 $AO_1,BO_2,CO_3$ 相交于一点 $P$．

\textbf{(3)} 由 (2) 知交点 $P$ 的坐标为 $\left(\frac{1}{4},\frac{1}{4},\frac{1}{4}\right)$．
\end{solution}

\answerproblem{二、}{15}
设 $N\geqslant 1$，$\mathcal{S}$ 是包含 $N$ 个整数的集合，满足如下的加性唯一性条件：

若 $n_k\in\mathcal{S}\ (1\leqslant k\leqslant 4)$ 且 $n_1+n_2=n_3+n_4$，则必有 $n_1=n_3$ 或 $n_1=n_4$．

令 $f(x):=\sum\limits_{n\in\mathcal{S}}\ee^{2\pi\ii nx}$．
\begin{enumerate}[label=(\arabic*)]
  \item 计算 $\displaystyle\int_0^1|f(x)|^2\dd x$，$\displaystyle\int_0^1|f(x)|^4\dd x$．
  \item 证明：$\displaystyle\int_0^1|f(x)|\dd x\geqslant\frac{1}{2\sqrt{N}}$．
\end{enumerate}
\begin{solution}
\textbf{(1)} 我们有
\begin{align*}
  |f(x)|^2=f(x)\overline{f(x)}
  &=\left(\sum_{m\in\mathcal{S}}\ee^{2\pi\ii mx}\right)\left(\sum_{n\in\mathcal{S}}\ee^{-2\pi\ii nx}\right)\\
  &\equiv N+g(x),
\end{align*}
其中，
\[
  g(x)=\sum_{\substack{m,n\in\mathcal{S}\\ m\neq n}}\ee^{2\pi\ii(m-n)x}.
\]
于是，注意到当整数 $k\neq 0$ 时，$\displaystyle\int_0^1\ee^{2\pi\ii kx}\dd x=0$，立即得到 $\displaystyle\int_0^1 g(x)\dd x=0$，从而
\[
  \int_0^1|f(x)|^2\dd x=N.
\]
而由题设条件，对于 $m_1\neq n_1$，$m_2\neq n_2$，若 $m_1\neq m_2$，则 $m_1-n_1\neq m_2-n_2$．于是，
\[
  g(x)=\sum_{k\in\mathcal{S}_g}\ee^{2\pi\ii kx},
\]
这里集合 $\mathcal{S}_g$ 由 $N^2-N$ 个两两不同的非零整数组成．这样，又有
\begin{align*}
  \int_0^1|f(x)|^4\dd x
  &=\int_0^1\left(N^2+2Ng(x)+|g(x)|^2\right)\dd x\\
  &=N^2+0+(N^2-N)=2N^2-N.
\end{align*}

\textbf{(2)} 对于本小题，我们提供两种方式的证明．

\textbf{法 I.} 注意到 $\displaystyle\int_0^1 g(x)\dd x=0$，我们有
\begin{align*}
  \sqrt{N}-\int_0^1|f(x)|\dd x
  &=\int_0^1\left(\sqrt{N}-\sqrt{N+g(x)}\right)\dd x
  =-\int_0^1\frac{g(x)}{\sqrt{N}+\sqrt{N+g(x)}}\dd x\\
  &=\int_0^1\left(\frac{g(x)}{2\sqrt{N}}-\frac{g(x)}{\sqrt{N}+\sqrt{N+g(x)}}\right)\dd x\\
  &=\int_0^1\frac{g^2(x)}{2\sqrt{N}\left(\sqrt{N}+\sqrt{N+g(x)}\right)^2}\dd x\\
  &\leqslant\int_0^1\frac{g^2(x)}{2N\sqrt{N}}\dd x
  =\frac{N^2-N}{2N\sqrt{N}}<\frac{\sqrt{N}}{2}.
\end{align*}
因此，
\[
  \int_0^1|f(x)|\dd x>\frac{\sqrt{N}}{2}.
\]

\textbf{法 II.} 由 Hölder 不等式，我们有
\begin{align*}
  N=\int_0^1|f(x)|^2\dd x
  &\leqslant\left(\int_0^1|f(x)|\dd x\right)^{\frac{2}{3}}
  \left(\int_0^1|f(x)|^4\dd x\right)^{\frac{1}{3}}\\
  &=\left(\int_0^1|f(x)|\dd x\right)^{\frac{2}{3}}\left(2N^2-N\right)^{\frac{1}{3}}.
\end{align*}
因此，
\[
  \int_0^1|f(x)|\dd x\geqslant\sqrt{\frac{N^3}{2N^2-N}}
  >\frac{\sqrt{N}}{\sqrt{2}}>\frac{\sqrt{N}}{2}.
\]
\end{solution}

\answerproblem{三、}{20}
设 $m,n$ 为大于 2 的整数，$a_1,\cdots,a_{m+1}$ 为任意 $m+1$ 个有理数，$\mathbb{Q}^{n\times n}$ 为有理数域上 $n$ 阶方阵全体．证明：
\begin{enumerate}[label=(\arabic*)]
  \item $\mathbb{Q}^{n\times n}$ 中存在 $m$ 个元素 $B_1,\cdots,B_m$ 使得行列式 $|B_j|=j\ (j=1,\cdots,m)$ 成立．
  \item $\mathbb{Q}^{n\times n}$ 中存在 $m$ 个元素 $A_1,\cdots,A_m$ 使得下列两条同时成立：
  \begin{enumerate}[label=(\roman*)]
    \item $|A_j|=a_j\ (j=1,\cdots,m)$；
    \item $|A_1-A_2-\cdots-A_m|=a_{m+1}$．
  \end{enumerate}
\end{enumerate}
\begin{solution}
\textbf{(1)} 令
\[
  B_1=\begin{pmatrix}1& & \\ & \ddots & \\ & & 1\end{pmatrix},\qquad
  B_2=\begin{pmatrix}1& & & \\ & 2 & & \\ & & \ddots & \\ & & & 1\end{pmatrix},\qquad
  \cdots,\qquad
  B_m=\begin{pmatrix}1& & & \\ & \ddots & & \\ & & 1 & \\ & & & m\end{pmatrix},
\]
则有
\[
  |B_j|=j,\qquad \forall\, j=1,\cdots,m.
\]

\textbf{(2)} 令
\[
  A_1=\begin{pmatrix}a_1 & 0 & & \\ 1 & 1 & & \\ & & \ddots & \\ & & & 1\end{pmatrix},\qquad
  A_2=\begin{pmatrix}a_2 & b & & \\ 0 & 1 & & \\ & & \ddots & \\ & & & 1\end{pmatrix},
\]
\[
  A_i=\begin{pmatrix}a_i & & & \\ & 1 & & \\ & & \ddots & \\ & & & 1\end{pmatrix},\qquad i=3,\cdots,m.
\]
则有
\begin{enumerate}[label=(\roman*)]
  \item $|A_j|=a_j$，$j=1,\cdots,m$；
  \item 
  \[
    |A_1-A_2-\cdots-A_m|=
    \begin{pmatrix}
      a_1-\sum\limits_{i=2}^{m}a_i & -b & & \\
      1 & -(m-2) & & \\
      & & \ddots & \\
      & & & -(m-2)
    \end{pmatrix}.
  \]
\end{enumerate}
于是有
\[
  |A_1-A_2-\cdots-A_m|=\left(-s(m-2)+b\right)[-(m-2)]^{m-2},
\]
其中
\[
  s=a_1-\sum_{i=2}^{m}a_i.
\]
取
\[
  b=s(m-2)+\frac{(-1)^ma_{m+1}}{(m-2)^{m-2}},
\]
则有 $b\in\mathbb{Q}$，从而 $A_1,\cdots,A_m$ 满足条件．证毕．
\end{solution}

\answerproblem{四、}{15}
设 $A$ 为 $n$ 阶方阵，其主对角线上的元素皆为 3，其余位置上的元素不是 2 就是 2029．证明：$\rank A=n$ 或 $n-1$．
\begin{solution}
令
\[
  B=\begin{pmatrix}2&\cdots&2\\ \vdots&\ddots&\vdots\\ 2&\cdots&2\end{pmatrix},
\]
且令 $C=A-B$，则有 $C$ 的主对角线上元素皆为 1，其余位置上的元素不是 0 就是 2027．

用 $\overline{m}$ 表示整数 $m$ 模 2027 所得的剩余类，将 $C$ 的每个元素模 2027，则有
\[
  \overline{C}=\begin{pmatrix}\overline{1}& & \overline{0}\\ & \ddots & \\ \overline{0} & & \overline{1}\end{pmatrix},
\]
于是
\[
  |\overline{C}|=\overline{1}=\overline{|C|}.
\]
从而 $A-B$ 是可逆的（即 $|A-B|\neq 0$）．于是
\[
  n=\rank(A-B)\leqslant \rank A+\rank B=\rank A+1,
\]
$\rank A=n-1$ 或 $n$．证毕．
\end{solution}

\answerproblem{五、}{15}
设 $k$ 为正整数．
\begin{enumerate}[label=(\arabic*)]
  \item 证明：对任何 $k\geqslant 1$，方程 $x^{3k+1}+x^2-6x+1=0$ 在 $\bigl(0,\frac{1}{2}\bigr)$ 内有唯一根 $x_k$；
  \item 证明：点列 $\{x_k\}$ 严格单减；
  \item 求极限 $\lim\limits_{k\to+\infty}x_k$．
\end{enumerate}
\begin{solution}
对于 $k\geqslant 1$，记 $f_k(x)=x^{3k+1}+x^2-6x+1$．

\textbf{(1)} 对于 $k\geqslant 1$，易见 $f_k$ 在 $\left[0,\frac{1}{2}\right]$ 上连续可微，$f_k(0)=1$，$f_k\left(\frac{1}{2}\right)<1+1-3+1=0$．由介值定理，存在 $x_k\in\left(0,\frac{1}{2}\right)$ 使得 $f_k(x_k)=0$．

另一方面，由二项展开可见 $2^{3k}=(1+1)^{3k}\geqslant 1+3k$，从而
\[
  f_k'(x)=(3k+1)x^{3k}+2x-6\leqslant\frac{3k+1}{2^{3k}}+1-6<0
  \quad \left(x\in\left[0,\frac{1}{2}\right]\right),
\]
因此，$f_k$ 在 $\left[0,\frac{1}{2}\right]$ 上严格单减．故 $f_k$ 在 $\left[0,\frac{1}{2}\right]$ 上至多只有一个零点．结合 $x_k$ 为 $f_k$ 的零点得到 $f_k$ 在 $\left(0,\frac{1}{2}\right)$ 内有唯一根 $x_k$．

\textbf{(2)} 对任何 $k\geqslant 1$，
\[
  f_{k+1}(x_k)=x_k^{3k+4}+x_k^2-6x_k+1<x_k^{3k+1}+x_k^2-6x_k+1=f_k(x_k)=0.
\]
因此，结合 (1) 的结论得到 $x_{k+1}<x_k$．即 $\{x_k\}$ 严格单减．

\textbf{(3)} 由于 $\{x_k\}$ 是严格单减的正数列，因此 $\lim\limits_{k\to+\infty}x_k$ 存在，设为 $\bar{x}$．则注意到 $0<x_k^{3k+1}<\frac{1}{2^{3k+1}}$，可得 $\lim\limits_{k\to+\infty}x_k^{3k+1}=0$．因此，
\[
  \bar{x}^2-6\bar{x}+1=0.
\]
解得 $\bar{x}=3-2\sqrt{2}$ 或 $\bar{x}=3+2\sqrt{2}$（舍去）．最后得到 $\lim\limits_{k\to+\infty}x_k=3-2\sqrt{2}$．
\end{solution}

\answerproblem{六、}{20}
设 $f\in C^1[0,1]$ 满足
\[
  f(x)\ln^2 x+xf'(x)\leqslant 0,\qquad \forall x\in(0,1).
\]
证明下列结论之一成立：
\begin{enumerate}[label=(\arabic*)]
  \item $\displaystyle\int_0^1 f(t)\dd t<\int_0^x f(t)\dd t\leqslant 0\ (\forall x\in(0,1))$；
  \item $f(x)=0\ (\forall x\in[0,1])$．
\end{enumerate}
\begin{solution}
考虑 $F(x):=\ee^{\frac{\ln^3 x}{3}}f(x)\ (x\in(0,1])$ 并补充定义 $F(0)=\lim\limits_{x\to0^+}F(x)=0$．则 $F$ 在 $[0,1]$ 上连续，在 $(0,1]$ 上可导．

由题设，
\[
  F'(x)=\left(\ee^{\frac{\ln^3 x}{3}}f(x)\right)'
  =\frac{1}{x}\ee^{\frac{\ln^3 x}{3}}\left(f(x)\ln^2 x+xf'(x)\right)\leqslant 0,
  \qquad \forall x\in(0,1).
\]
因此，$F$ 在 $[0,1]$ 上单调下降．

特别，$F(x)\leqslant F(0)=0\ (x\in[0,1])$．进而
\[
  f(x)\leqslant 0,\qquad \forall x\in[0,1].
\]
因此，
\[
  \int_0^1 f(t)\dd t=\int_0^x f(t)\dd t+\int_x^1 f(t)\dd t\leqslant\int_0^x f(t)\dd t,
  \qquad \forall x\in[0,1].
\]
若对每个 $x\in(0,1)$，$f(x)=0$，则由连续性，$f$ 在 $[0,1]$ 上恒为零，即情形 (2) 发生．

否则，存在一个 $\xi\in(0,1)$ 使得 $f(\xi)<0$．从而 $F(\xi)<0$．进而对任何 $t\in[\xi,1]$，$F(t)\leqslant F(\xi)<0$，从而又有 $f(t)<0$．于是对任何 $x\in(0,1)$，$\displaystyle\int_x^1 f(t)\dd t<0$．所以
\[
  \int_0^1 f(t)\dd t<\int_0^x f(t)\dd t,\qquad \forall x\in(0,1),
\]
即情形 (1) 发生．$\square$
\end{solution}
